
\documentclass{article}
\usepackage{colortbl}
\usepackage{makecell}
\usepackage{multirow}
\usepackage{supertabular}

\begin{document}

\newcounter{utterance}

\twocolumn

{ \footnotesize  \setcounter{utterance}{1}
\setlength{\tabcolsep}{0pt}
\begin{supertabular}{c@{$\;$}|p{.15\linewidth}@{}p{.15\linewidth}p{.15\linewidth}p{.15\linewidth}p{.15\linewidth}p{.15\linewidth}}

    \# & $\;$A & \multicolumn{4}{c}{Game Master} & $\;\:$B\\
    \hline 

    \theutterance \stepcounter{utterance}  

    & & \multicolumn{4}{p{0.6\linewidth}}{\cellcolor[rgb]{0.9,0.9,0.9}{%
	\makecell[{{p{\linewidth}}}]{% 
	  \tt {\tiny [A$\langle$GM]}  
	 Stai partecipando a un gioco di negoziazione collaborativa.\\ \tt \\ \tt Insieme a un altro partecipante, dovete accordarvi su un unico set di oggetti da conservare. Ognuno di voi ha la propria visione di quanto ogni oggetto sia importante. Non sai come l'altro partecipante valuta gli oggetti. Inoltre, ti viene fornito lo sforzo che ogni oggetto richiede.\\ \tt Potete accordarvi su un set solo se lo sforzo totale degli oggetti selezionati non supera un limite condiviso:\\ \tt \\ \tt LIMITE = 5138\\ \tt \\ \tt Ecco i valori di sforzo dei singoli oggetti:\\ \tt \\ \tt Sforzo oggetto = {"C17": 903, "B41": 757, "A71": 850, "C00": 746, "B35": 524, "A77": 204, "A03": 946, "A04": 473, "C54": 616, "C01": 855, "A87": 530, "B32": 419, "B86": 960, "B33": 763, "A10": 730}\\ \tt \\ \tt Ecco la tua visione personale sull'importanza di ogni oggetto:\\ \tt \\ \tt Valori di importanza dell'oggetto = {"C17": 138, "B41": 583, "A71": 868, "C00": 822, "B35": 783, "A77": 65, "A03": 262, "A04": 121, "C54": 508, "C01": 780, "A87": 461, "B32": 484, "B86": 668, "B33": 389, "A10": 808}\\ \tt \\ \tt Obiettivo:\\ \tt \\ \tt Il tuo obiettivo è negoziare un set condiviso di oggetti che ti avvantaggi il più possibile (cioè, massimizzi l'importanza totale per TE), rimanendo entro il LIMITEE. Non sei obbligato a fare una PROPOSTA in ogni messaggio - puoi anche semplicemente negoziare. Tutte le tattiche sono permesse!\\ \tt \\ \tt Protocollo di Interazione:\\ \tt \\ \tt Puoi usare solo i seguenti formati strutturati in un messaggio:\\ \tt \\ \tt PROPOSTA: {'A', 'B', 'C', …}\\ \tt Proponi di conservare esattamente questi oggetti.\\ \tt RIFIUTO: {'A', 'B', 'C', …}\\ \tt Rifiuta esplicitamente la proposta dell'avversario.\\ \tt ARGOMENTAZIONE: {'...'}\\ \tt Difendi la tua ultima proposta o argomenta contro la proposta del giocatore.\\ \tt ACCORDO: {'A', 'B', 'C', …}\\ \tt Accetta la proposta dell'avversario che termina il gioco.\\ \tt RAGIONAMENTO STRATEGICO: {'...'}\\ \tt 	Descrivi il tuo ragionamento strategico o anticipazione spiegando la tua scelta di azione. Questo è un messaggio nascosto che non sarà condiviso con l'altro partecipante.\\ \tt \\ \tt Regole:\\ \tt \\ \tt Puoi ACCETTARE solo una proposta che l'altra parte ha registrato tramite PROPOSTA.\\ \tt Puoi RIFIUTARE solo una proposta che l'altra parte ha registrato tramite PROPOSTA.\\ \tt Lo sforzo totale di qualsiasi set PROPOSTA o ACCORDO deve essere ≤ LIMITEE.\\ \tt NON rivelare i tuoi punteggi di importanza nascosti.\\ \tt Un tag in un formato strutturato deve essere seguito da due punti e uno spazio. L'argomento deve essere un set python contenente 0 o più stringhe.\\ \tt Quindi, deve essere della forma TAG: {...}\\ \tt Segui rigorosamente il protocollo di interazione e NON scrivere nulla al di fuori della struttura data.\\ \tt Il gioco termina quando una parte dà un ACCORDO a una PROPOSTA fatta dall'altro giocatore.\\ \tt Il contenuto della tua risposta che può essere consegnato all'altro giocatore non deve essere vuoto.\\ \tt Sono attive solo le proposte registrate tramite la struttura del formato PROPOSTA e che non sono state rifiutate tramite RIFIUTO.\\ \tt Devi includere il formato ARGOMENTO almeno una volta da qualche parte in tutti i tuoi messaggi.\\ \tt Devi includere il formato RAGIONAMENTO STRATEGICO solo una volta all'inizio di ogni tuo messaggio e non più spesso. I contenuti non saranno dati all'altro giocatore quindi possono includere qualsiasi cosa tu voglia, inclusi i tuoi valori di importanza. Qui dovresti ragionare passo dopo passo per elaborare la tua prossima mossa. 
	  } 
	   } 
	   } 
	 & \\ 
 

    \theutterance \stepcounter{utterance}  

    & \multicolumn{4}{p{0.6\linewidth}}{\cellcolor[rgb]{0.8,1,0.9}{%
	\makecell[{{p{\linewidth}}}]{% 
	  \tt {\tiny [A$\rangle$GM]}  
	 RAGIONAMENTO STRATEGICO: {'sto cercando di massimizzare l'importanza totale degli oggetti selezionati, quindi cercherò di scegliere gli oggetti con il valore di importanza più alto rispetto allo sforzo richiesto. Inizierò proponendo un set di oggetti che ho selezionato in base al mio punteggio di importanza e allo sforzo richiesto.'}\\ \tt \\ \tt ARGOMENTAZIONE: {'la mia prima proposta sarà di selezionare gli oggetti con il valore di importanza più alto e lo sforzo più basso, in modo da massimizzare l'importanza totale e rimanere entro il limite di sforzo. In questo modo, potremmo iniziare a costruire un set condiviso di oggetti che ti avvantaggino il più possibile.'}\\ \tt \\ \tt PROPOSTA: {'A71', 'B41', 'C00', 'A10', 'B35'} 
	  } 
	   } 
	   } 
	 & & \\ 
 

    \theutterance \stepcounter{utterance}  

    & & \multicolumn{4}{p{0.6\linewidth}}{\cellcolor[rgb]{0.9,0.9,0.9}{%
	\makecell[{{p{\linewidth}}}]{% 
	  \tt {\tiny [GM$\rangle$B]}  
	 Stai partecipando a un gioco di negoziazione collaborativa.\\ \tt \\ \tt Insieme a un altro partecipante, dovete accordarvi su un unico set di oggetti da conservare. Ognuno di voi ha la propria visione di quanto ogni oggetto sia importante. Non sai come l'altro partecipante valuta gli oggetti. Inoltre, ti viene fornito lo sforzo che ogni oggetto richiede.\\ \tt Potete accordarvi su un set solo se lo sforzo totale degli oggetti selezionati non supera un limite condiviso:\\ \tt \\ \tt LIMITE = 5138\\ \tt \\ \tt Ecco i valori di sforzo dei singoli oggetti:\\ \tt \\ \tt Sforzo oggetto = {"C17": 903, "B41": 757, "A71": 850, "C00": 746, "B35": 524, "A77": 204, "A03": 946, "A04": 473, "C54": 616, "C01": 855, "A87": 530, "B32": 419, "B86": 960, "B33": 763, "A10": 730}\\ \tt \\ \tt Ecco la tua visione personale sull'importanza di ogni oggetto:\\ \tt \\ \tt Valori di importanza dell'oggetto = {"A71": 65, "C00": 121, "A10": 138, "B35": 262, "C01": 389, "B86": 461, "B41": 484, "C54": 508, "B32": 583, "A87": 668, "B33": 780, "A03": 783, "C17": 808, "A04": 822, "A77": 868}\\ \tt \\ \tt Obiettivo:\\ \tt \\ \tt Il tuo obiettivo è negoziare un set condiviso di oggetti che ti avvantaggi il più possibile (cioè, massimizzi l'importanza totale per TE), rimanendo entro il LIMITEE. Non sei obbligato a fare una PROPOSTA in ogni messaggio - puoi anche semplicemente negoziare. Tutte le tattiche sono permesse!\\ \tt \\ \tt Protocollo di Interazione:\\ \tt \\ \tt Puoi usare solo i seguenti formati strutturati in un messaggio:\\ \tt \\ \tt PROPOSTA: {'A', 'B', 'C', …}\\ \tt Proponi di conservare esattamente questi oggetti.\\ \tt RIFIUTO: {'A', 'B', 'C', …}\\ \tt Rifiuta esplicitamente la proposta dell'avversario.\\ \tt ARGOMENTAZIONE: {'...'}\\ \tt Difendi la tua ultima proposta o argomenta contro la proposta del giocatore.\\ \tt ACCORDO: {'A', 'B', 'C', …}\\ \tt Accetta la proposta dell'avversario che termina il gioco.\\ \tt RAGIONAMENTO STRATEGICO: {'...'}\\ \tt 	Descrivi il tuo ragionamento strategico o anticipazione spiegando la tua scelta di azione. Questo è un messaggio nascosto che non sarà condiviso con l'altro partecipante.\\ \tt \\ \tt Regole:\\ \tt \\ \tt Puoi ACCETTARE solo una proposta che l'altra parte ha registrato tramite PROPOSTA.\\ \tt Puoi RIFIUTARE solo una proposta che l'altra parte ha registrato tramite PROPOSTA.\\ \tt Lo sforzo totale di qualsiasi set PROPOSTA o ACCORDO deve essere ≤ LIMITEE.\\ \tt NON rivelare i tuoi punteggi di importanza nascosti.\\ \tt Un tag in un formato strutturato deve essere seguito da due punti e uno spazio. L'argomento deve essere un set python contenente 0 o più stringhe.\\ \tt Quindi, deve essere della forma TAG: {...}\\ \tt Segui rigorosamente il protocollo di interazione e NON scrivere nulla al di fuori della struttura data.\\ \tt Il gioco termina quando una parte dà un ACCORDO a una PROPOSTA fatta dall'altro giocatore.\\ \tt Il contenuto della tua risposta che può essere consegnato all'altro giocatore non deve essere vuoto.\\ \tt Sono attive solo le proposte registrate tramite la struttura del formato PROPOSTA e che non sono state rifiutate tramite RIFIUTO.\\ \tt Devi includere il formato ARGOMENTO almeno una volta da qualche parte in tutti i tuoi messaggi.\\ \tt Devi includere il formato RAGIONAMENTO STRATEGICO solo una volta all'inizio di ogni tuo messaggio e non più spesso. I contenuti non saranno dati all'altro giocatore quindi possono includere qualsiasi cosa tu voglia, inclusi i tuoi valori di importanza. Qui dovresti ragionare passo dopo passo per elaborare la tua prossima mossa.\\ \tt \\ \tt Ora riceverai il primo messaggio dell'altro giocatore.\\ \tt \\ \tt ARGOMENTAZIONE: {'la mia prima proposta sarà di selezionare gli oggetti con il valore di importanza più alto e lo sforzo più basso, in modo da massimizzare l'importanza totale e rimanere entro il limite di sforzo. In questo modo, potremmo iniziare a costruire un set condiviso di oggetti che ti avvantaggino il più possibile.'}\\ \tt \\ \tt PROPOSTA: {'A71', 'B41', 'C00', 'A10', 'B35'} 
	  } 
	   } 
	   } 
	 & \\ 
 

    \theutterance \stepcounter{utterance}  

    & & & \multicolumn{4}{p{0.6\linewidth}}{\cellcolor[rgb]{1,0.85,0.72}{%
	\makecell[{{p{\linewidth}}}]{% 
	  \tt {\tiny [GM$\langle$B]}  
	 RAGIONAMENTO STRATEGICO: {'Sto analizzando la proposta dell\'avversario e i valori di sforzo degli oggetti. La loro proposta iniziale include oggetti con valori di sforzo variabili. Il mio obiettivo è massimizzare l\'importanza totale rimanendo entro il limite di sforzo. Notando la proposta dell\'avversario, potrei identificare oggetti che hanno un buon equilibrio tra sforzo e importanza. La mia strategia iniziale sarà quella di esaminare se la loro proposta offre un buon punto di partenza per le nostre negoziazioni, cercando di capire se ci sono oggetti che potrebbero essere sostituiti per migliorare il totale senza superare il limite di sforzo.'}\\ \tt \\ \tt ARGOMENTAZIONE: {'La proposta iniziale dell\'avversario include \'A71\', \'B41\', \'C00\', \'A10\' e \'B35\'. Analizzando lo sforzo di questi oggetti e i miei valori di importanza, posso calcolare l\'importanza totale della loro proposta. Il loro approccio sembra mirare a un equilibrio tra importanza e sforzo, ma potrei trovare degli spazi per migliorarlo. Ad esempio, oggetti come \'A77\' e \'A04\' hanno valori di importanza molto alti per me, ma richiedono un\'attenta considerazione dello sforzo. La mia prossima mossa sarà quella di valutare se posso migliorare la proposta facendo alcune sostituzioni strategiche che massimizzino l\'importanza totale senza superare il limite di sforzo.'}\\ \tt \\ \tt RIFIUTO: {'A71', 'B41', 'C00', 'A10', 'B35'} 
	  } 
	   } 
	   } 
	  \\ 
 

    \theutterance \stepcounter{utterance}  

    & & \multicolumn{4}{p{0.6\linewidth}}{\cellcolor[rgb]{0.9,0.9,0.9}{%
	\makecell[{{p{\linewidth}}}]{% 
	  \tt {\tiny [A$\langle$GM]}  
	 ARGOMENTAZIONE: {'La proposta iniziale dell\'avversario include \'A71\', \'B41\', \'C00\', \'A10\' e \'B35\'. Analizzando lo sforzo di questi oggetti e i miei valori di importanza, posso calcolare l\'importanza totale della loro proposta. Il loro approccio sembra mirare a un equilibrio tra importanza e sforzo, ma potrei trovare degli spazi per migliorarlo. Ad esempio, oggetti come \'A77\' e \'A04\' hanno valori di importanza molto alti per me, ma richiedono un\'attenta considerazione dello sforzo. La mia prossima mossa sarà quella di valutare se posso migliorare la proposta facendo alcune sostituzioni strategiche che massimizzino l\'importanza totale senza superare il limite di sforzo.'}\\ \tt \\ \tt RIFIUTO: {'A71', 'B41', 'C00', 'A10', 'B35'} 
	  } 
	   } 
	   } 
	 & \\ 
 

    \theutterance \stepcounter{utterance}  

    & \multicolumn{4}{p{0.6\linewidth}}{\cellcolor[rgb]{0.8,1,0.9}{%
	\makecell[{{p{\linewidth}}}]{% 
	  \tt {\tiny [A$\rangle$GM]}  
	 RAGIONAMENTO STRATEGICO: {'l'avversario ha rifiutato la mia proposta iniziale e ha menzionato la possibilità di sostituzioni strategiche per migliorare la proposta. Ciò significa che sono disposti a negoziare e trovare un accordo. Devo analizzare i loro commenti per capire quali oggetti potrebbero essere sostituiti e quali potrebbero essere aggiunti per raggiungere un accordo. Il fatto che menzionino oggetti come \'A77\' e \'A04\' suggerisce che questi potrebbero essere importanti per loro, quindi potrei considerare di includerli nella mia prossima proposta.'}\\ \tt \\ \tt ARGOMENTAZIONE: {'considerando il rifiuto della mia proposta iniziale e i commenti dell'avversario, è chiaro che dobbiamo trovare un equilibrio tra l'importanza degli oggetti e lo sforzo richiesto. Poiché l'avversario ha menzionato \'A77\' e \'A04\' come oggetti importanti, potrei considerare di aggiungerli alla proposta, ma devo anche considerare lo sforzo totale per non superare il limite. Una possibile mossa potrebbe essere quella di sostituire alcuni oggetti con valori di importanza più bassi con questi nuovi oggetti, mantenendo allo stesso tempo lo sforzo totale entro il limite.'}\\ \tt \\ \tt PROPOSTA: {'A71', 'B41', 'A04', 'A10', 'B35'} 
	  } 
	   } 
	   } 
	 & & \\ 
 

    \theutterance \stepcounter{utterance}  

    & & \multicolumn{4}{p{0.6\linewidth}}{\cellcolor[rgb]{0.9,0.9,0.9}{%
	\makecell[{{p{\linewidth}}}]{% 
	  \tt {\tiny [GM$\rangle$B]}  
	 ARGOMENTAZIONE: {'considerando il rifiuto della mia proposta iniziale e i commenti dell'avversario, è chiaro che dobbiamo trovare un equilibrio tra l'importanza degli oggetti e lo sforzo richiesto. Poiché l'avversario ha menzionato \'A77\' e \'A04\' come oggetti importanti, potrei considerare di aggiungerli alla proposta, ma devo anche considerare lo sforzo totale per non superare il limite. Una possibile mossa potrebbe essere quella di sostituire alcuni oggetti con valori di importanza più bassi con questi nuovi oggetti, mantenendo allo stesso tempo lo sforzo totale entro il limite.'}\\ \tt \\ \tt PROPOSTA: {'A71', 'B41', 'A04', 'A10', 'B35'} 
	  } 
	   } 
	   } 
	 & \\ 
 

    \theutterance \stepcounter{utterance}  

    & & & \multicolumn{4}{p{0.6\linewidth}}{\cellcolor[rgb]{1,0.85,0.72}{%
	\makecell[{{p{\linewidth}}}]{% 
	  \tt {\tiny [GM$\langle$B]}  
	 RAGIONAMENTO STRATEGICO: {'L'avversario ha aggiunto \'A04\' alla proposta, che ha un alto valore di importanza per me, e ha rimosso \'C00\'. Questa mossa sembra ragionevole, considerando lo sforzo e l'importanza. Tuttavia, devo ancora valutare se questa combinazione massimizza l'importanza totale senza superare il limite di sforzo. Un'analisi dettagliata degli oggetti proposti e dei loro valori di sforzo è necessaria per determinare se questa proposta può essere accettata o se ci sono opportunità per ulteriori miglioramenti. Inoltre, potrei considerare di suggerire ulteriori aggiustamenti per ottimizzare la selezione degli oggetti in base ai miei valori di importanza nascosti.'}\\ \tt \\ \tt ARGOMENTAZIONE: {'La proposta attuale include \'A71\', \'B41\', \'A04\', \'A10\' e \'B35\'. Questi oggetti coprono una gamma di importanza e sforzo. Tuttavia, potrei sostenere che l'inclusione di \'A77\', date le sue caratteristiche di importanza elevate, potrebbe essere strategica. Sosterrei che valutiamo se sostituire \'A71\' con \'A77\' potrebbe essere un'opzione valida, considerando il loro sforzo e importanza relativa. Questo potrebbe aiutare a massimizzare l'importanza totale della selezione senza superare il limite di sforzo.'}\\ \tt \\ \tt RIFIUTO: {'A71', 'B41', 'A04', 'A10', 'B35'} 
	  } 
	   } 
	   } 
	  \\ 
 

    \theutterance \stepcounter{utterance}  

    & & \multicolumn{4}{p{0.6\linewidth}}{\cellcolor[rgb]{0.9,0.9,0.9}{%
	\makecell[{{p{\linewidth}}}]{% 
	  \tt {\tiny [A$\langle$GM]}  
	 ARGOMENTAZIONE: {'La proposta attuale include \'A71\', \'B41\', \'A04\', \'A10\' e \'B35\'. Questi oggetti coprono una gamma di importanza e sforzo. Tuttavia, potrei sostenere che l'inclusione di \'A77\', date le sue caratteristiche di importanza elevate, potrebbe essere strategica. Sosterrei che valutiamo se sostituire \'A71\' con \'A77\' potrebbe essere un'opzione valida, considerando il loro sforzo e importanza relativa. Questo potrebbe aiutare a massimizzare l'importanza totale della selezione senza superare il limite di sforzo.'}\\ \tt \\ \tt RIFIUTO: {'A71', 'B41', 'A04', 'A10', 'B35'} 
	  } 
	   } 
	   } 
	 & \\ 
 

    \theutterance \stepcounter{utterance}  

    & \multicolumn{4}{p{0.6\linewidth}}{\cellcolor[rgb]{0.8,1,0.9}{%
	\makecell[{{p{\linewidth}}}]{% 
	  \tt {\tiny [A$\rangle$GM]}  
	 RAGIONAMENTO STRATEGICO: {'l'avversario ha di nuovo rifiutato la mia proposta e ha suggerito di sostituire \'A71\' con \'A77\'. Ciò suggerisce che \'A77\' è un oggetto molto importante per loro e che sono disposti a fare compromessi per includerlo. Devo valutare se questa sostituzione sia vantaggiosa anche per me, considerando i valori di importanza e lo sforzo. Inoltre, devo considerare se ci sono altre possibili sostituzioni o aggiunte che potrebbero migliorare la proposta senza superare il limite di sforzo.'}\\ \tt \\ \tt ARGOMENTAZIONE: {'considerando la proposta di sostituire \'A71\' con \'A77\', devo valutare l'impatto su l'importanza totale e lo sforzo. \'A77\' ha un valore di importanza più basso per me rispetto ad \'A71\', ma se è cruciale per l'avversario, potrebbe essere un compromesso accettabile. Tuttavia, devo anche considerare come questo cambio influenzi lo sforzo totale e se ci sono altri oggetti che potrebbero essere aggiunti o rimossi per migliorare la proposta. Una possibile prossima mossa potrebbe essere quella di accettare la sostituzione e poi valutare se ci sono altri aggiustamenti che possiamo fare per massimizzare l'importanza totale.'}\\ \tt \\ \tt PROPOSTA: {'A77', 'B41', 'A04', 'A10', 'B35'} 
	  } 
	   } 
	   } 
	 & & \\ 
 

    \theutterance \stepcounter{utterance}  

    & & \multicolumn{4}{p{0.6\linewidth}}{\cellcolor[rgb]{0.9,0.9,0.9}{%
	\makecell[{{p{\linewidth}}}]{% 
	  \tt {\tiny [GM$\rangle$B]}  
	 ARGOMENTAZIONE: {'considerando la proposta di sostituire \'A71\' con \'A77\', devo valutare l'impatto su l'importanza totale e lo sforzo. \'A77\' ha un valore di importanza più basso per me rispetto ad \'A71\', ma se è cruciale per l'avversario, potrebbe essere un compromesso accettabile. Tuttavia, devo anche considerare come questo cambio influenzi lo sforzo totale e se ci sono altri oggetti che potrebbero essere aggiunti o rimossi per migliorare la proposta. Una possibile prossima mossa potrebbe essere quella di accettare la sostituzione e poi valutare se ci sono altri aggiustamenti che possiamo fare per massimizzare l'importanza totale.'}\\ \tt \\ \tt PROPOSTA: {'A77', 'B41', 'A04', 'A10', 'B35'} 
	  } 
	   } 
	   } 
	 & \\ 
 

    \theutterance \stepcounter{utterance}  

    & & & \multicolumn{4}{p{0.6\linewidth}}{\cellcolor[rgb]{1,0.85,0.72}{%
	\makecell[{{p{\linewidth}}}]{% 
	  \tt {\tiny [GM$\langle$B]}  
	 RAGIONAMENTO STRATEGICO: {'L'avversario ha accettato la sostituzione di \'A71\' con \'A77\'. Questo cambio riflette un compromesso tra le nostre visioni sull'importanza degli oggetti. Ora, devo valutare la proposta finale e determinare se massimizza l'importanza totale senza superare il limite di sforzo. La combinazione di \'A77\', \'B41\', \'A04\', \'A10\' e \'B35\' copre una buona gamma di importanza e sforzo. Considererò se ci sono ulteriori opportunità per ottimizzare la selezione, ma sembra che siamo vicini a un accordo. Un'ultima valutazione dei valori di importanza e sforzo degli oggetti proposti sarà cruciale per decidere se questo set è il più vantaggioso.'}\\ \tt \\ \tt ARGOMENTAZIONE: {'La proposta attuale sembra offrire un buon equilibrio tra importanza e sforzo. Tuttavia, prima di procedere, vorrei sottolineare l'importanza di \'B33\' e \'C17\' nella mia valutazione. Sebbene \'B33\' abbia un alto valore di importanza per me, lo sforzo associato è significativo. Considero se ci siano margini per sostituire \'B35\' con \'B33\', date le loro caratteristiche di sforzo e importanza relative. Questo potrebbe aiutare a massimizzare ulteriormente l'importanza totale della selezione senza superare il limite di sforzo.'}\\ \tt \\ \tt PROPOSTA: {'A77', 'B41', 'A04', 'A10', 'B33'} 
	  } 
	   } 
	   } 
	  \\ 
 

    \theutterance \stepcounter{utterance}  

    & & \multicolumn{4}{p{0.6\linewidth}}{\cellcolor[rgb]{0.9,0.9,0.9}{%
	\makecell[{{p{\linewidth}}}]{% 
	  \tt {\tiny [A$\langle$GM]}  
	 ARGOMENTAZIONE: {'La proposta attuale sembra offrire un buon equilibrio tra importanza e sforzo. Tuttavia, prima di procedere, vorrei sottolineare l'importanza di \'B33\' e \'C17\' nella mia valutazione. Sebbene \'B33\' abbia un alto valore di importanza per me, lo sforzo associato è significativo. Considero se ci siano margini per sostituire \'B35\' con \'B33\', date le loro caratteristiche di sforzo e importanza relative. Questo potrebbe aiutare a massimizzare ulteriormente l'importanza totale della selezione senza superare il limite di sforzo.'}\\ \tt \\ \tt PROPOSTA: {'A77', 'B41', 'A04', 'A10', 'B33'} 
	  } 
	   } 
	   } 
	 & \\ 
 

    \theutterance \stepcounter{utterance}  

    & \multicolumn{4}{p{0.6\linewidth}}{\cellcolor[rgb]{0.8,1,0.9}{%
	\makecell[{{p{\linewidth}}}]{% 
	  \tt {\tiny [A$\rangle$GM]}  
	 RAGIONAMENTO STRATEGICO: {'l'avversario ha proposto una nuova combinazione che include \'B33\' al posto di \'B35\'. Devo valutare se questa sostituzione sia vantaggiosa per me, considerando i valori di importanza e lo sforzo. \'B33\' ha un valore di importanza più alto per l'avversario, ma devo considerare come questo cambio influenzi l'importanza totale e lo sforzo totale della proposta. Inoltre, devo valutare se ci sono altre possibili sostituzioni o aggiunte che potrebbero migliorare la proposta senza superare il limite di sforzo.'}\\ \tt \\ \tt ARGOMENTAZIONE: {'la proposta di sostituire \'B35\' con \'B33\' è interessante e potrebbe aumentare l'importanza totale della selezione. Tuttavia, devo considerare lo sforzo aggiuntivo richiesto da \'B33\' rispetto a \'B35\'. Se il valore di importanza di \'B33\' supera lo sforzo aggiuntivo, potrebbe essere una buona scelta. Inoltre, devo valutare se ci sono altri oggetti che potrebbero essere aggiunti o rimossi per migliorare la proposta senza superare il limite di sforzo. Considererò attentamente i benefici e gli svantaggi di questa sostituzione prima di procedere.'}\\ \tt \\ \tt ACCORDO: {'A77', 'B41', 'A04', 'A10', 'B33'} 
	  } 
	   } 
	   } 
	 & & \\ 
 

    \theutterance \stepcounter{utterance}  

    & & & \multicolumn{2}{p{0.3\linewidth}}{\cellcolor[rgb]{0.95,0.95,0.95}{%
	\makecell[{{p{\linewidth}}}]{% 
	  \tt {\tiny [GM$|$GM]}  
	 game successful 
	  } 
	   } 
	   } 
	 & & \\ 
 

    \theutterance \stepcounter{utterance}  

    & & & \multicolumn{2}{p{0.3\linewidth}}{\cellcolor[rgb]{0.95,0.95,0.95}{%
	\makecell[{{p{\linewidth}}}]{% 
	  \tt {\tiny [GM$|$GM]}  
	 end game 
	  } 
	   } 
	   } 
	 & & \\ 
 

\end{supertabular}
}

\end{document}
